\chapter[Estado del arte]{Estado del arte}
\label{cp:estado_arte}
% Extensión estimada: 12--15 páginas.
% Fuentes: Docs/estado_del_arte/resumen_papers.tex (fichas de ~40 trabajos y tabla comparativa de
% arquitecturas), sección "Estado del Arte" de la Propuesta, antecedentes del Plan de trabajo y
% Docs/tabla_resultados_sota (tablas comparativas de resultados).

\pendiente{Este capítulo revisa los trabajos relacionados, organizados por tema. Cada sección cierra
indicando qué toma este trabajo de esa línea o en qué se diferencia.}

\section{Inteligencia artificial para el diagnóstico de lesiones de piel}
\label{sec:ea_clasificacion}
\pendiente{Se recorre la evolución desde las CNNs clásicas hacia los transformers, con resultados de
referencia en ISIC y HAM10000, revisiones sistemáticas y meta-análisis sobre melanoma y carcinoma
basocelular, y los problemas recurrentes del área: desbalance, generalización y sesgo por fototipo.}

\section{Modelos fundacionales y visión--lenguaje en dermatología}
\label{sec:ea_fundacionales}
\pendiente{Se presentan PanDerm, MONET, Derm1M y DermLIP, y MAKE: modelos pre-entrenados a gran escala con
texto clínico y su desempeño \emph{zero-shot}. Esta revisión justifica la elección de DermLIP.}

\section{Modelos visión--lenguaje médicos generales}
\label{sec:ea_vlm_generales}
\pendiente{Se describen LLaVA-Med, Med-Flamingo, PeFoMed y MedImageInsight, su cobertura de modalidades y la
evidencia disponible sobre su desempeño en dermatología.}

\section{Generación de reportes dermatológicos}
\label{sec:ea_reportes}
\pendiente{Se revisan SkinGPT-4, SkinVL (MM-Skin), DermatoLlama/SCALEMED, la destilación de conocimiento de
Chakrabarti et al., Skin-R1, SkinGEN y ensambles con LLM para redactar reportes, junto con la forma en que
cada uno evalúa (métricas con referencia, evaluación por especialistas) y el hardware que requiere.}

\section{Conjuntos de datos}
\label{sec:ea_datasets}
\pendiente{Se comparan HAM10000, ISIC, PH2, SkinCaRe (SkinCAP y SkinCoT), MM-Skin, Derm1M, HIBA,
Fitzpatrick17k y SkinCon en modalidad, tamaño, etiquetas, disponibilidad de texto y limitaciones conocidas,
incluida la ausencia de identificador de paciente en HAM10000 y la presencia de varias imágenes por lesión.}

\section{Sesgo, equidad y generalización}
\label{sec:ea_sesgo}
\pendiente{Se revisan las disparidades de desempeño por tono de piel (Groh et al., Daneshjou et al.) y la
evaluación en poblaciones distintas a la de entrenamiento, que conecta con la evaluación en HIBA.}

\section{Síntesis y brecha que aborda esta tesina}
\label{sec:ea_brecha}
\pendiente{Una tabla comparativa resume los trabajos (modelo, tarea, datos, métricas, hardware). La brecha
identificada es triple: faltan comparaciones bajo un mismo protocolo entre clasificadores supervisados, VLMs
generales y modelos visión--lenguaje dermatológicos; la mayoría de los generadores se evalúa por similitud
textual y no por fidelidad diagnóstica; y pocos trabajos consideran la incertidumbre del diagnóstico o el
hardware accesible.}
